%%%%%%%%%%  scan idx 81  |  folio 70  %%%%%%%%%%

\section*{3. Los haces coherentes $\underline D_X^i$}

Un funtor importante que por lo común no es prorepresentable, pero que
sin embargo admite un objeto formal versal, es el funtor de deformaciones.
Con este propósito, recordamos algunos hechos básicos sobre las extensiones
de anillos conmutativos y su teoría de obstrucciones. Esta sección no
contiene nada nuevo, salvo quizá por su uso sistemático de la cohomología de
Grothendieck sobre un sitio. Su finalidad es hacer esta exposición tan
autocontenida como sea posible; de hecho, demostramos algo más de lo que se
necesitará después. Los resultados básicos se enuncian en el Teorema 3.12.
En toda la sección, todos los anillos son conmutativos y poseen unidad.

%%%%%%%%%%  scan idx 82  |  folio 71  %%%%%%%%%%

Dado un anillo $R$, sea $\underline C_R$ la categoría de las
$R$-álgebras. Para todo $A\in\operatorname{Ob}(\underline C_R)$,
pongamos
\[
 \operatorname{Cov}(A)=
 \{\text{flechas suprayectivas }\varphi:A'\longrightarrow A
   \text{ en }\underline C_R\}.
\]
Esto define una pretopología, y por tanto un sitio, sobre
$\underline C_R$. Para todo $A\in\operatorname{Ob}(\underline C_R)$,
sea $\underline C_{A|R}$ la categoría de las flechas de
$\underline C_R$ con destino $A$. Dotamos a $\underline C_{A|R}$ de la
topología inducida por el funtor de olvido
\[
 j_A:\underline C_{A|R}\longrightarrow\underline C_R.
\]
La categoría $\underline C_{A|R}$ tiene un objeto final
$(1_A:A\longrightarrow A)$ y un objeto inicial
$(i_A:R\longrightarrow A)$, y admite tanto productos fibrados como
coproductos. Cuando no haya riesgo de confusión, escribiremos simplemente
$B$ para un objeto $(B\longrightarrow A)$ de $\underline C_{A|R}$.

Sea $\mathscr C$ una categoría abeliana. Un haz con valores en $\mathscr C$
sobre $\underline C_{A|R}$ es, por definición, un funtor
\[
 F:\underline C_{A|R}^{\mathrm{op}}\longrightarrow\mathscr C
\]
tal que, para toda flecha suprayectiva $B'\longrightarrow B$ en
$\underline C_{A|R}$, la sucesión
\[
 0\longrightarrow F(B)\longrightarrow F(B')
 \mathrel{\substack{\longrightarrow\\[-.6ex]\longrightarrow}}
 F(B'\mathbin{\times_B}B')
\]
es exacta. Los prehaces \underline{aditivos} son los que conservan los
coproductos; así,
\[
 F(B_1\otimes_R B_2)=F(B_1)\oplus F(B_2).
\]
Por ejemplo, un $A$-módulo $M$ define el funtor
\[
 \operatorname{Der}_R(-,M):
 \underline C_{A|R}^{\mathrm{op}}\longrightarrow(A\text{-mod}),
\]
que envía $B$ a las derivaciones $R$-lineales
$\operatorname{Der}_R(B,M)$, considerando $M$ como $B$-módulo por medio
de la aplicación estructural $B\longrightarrow A$. Este es claramente un
haz aditivo sobre $\underline C_{A|R}$. Aquí $(A\text{-mod})$ denota la
categoría de los $A$-módulos.

Se dice que un prehaz $F$ sobre $\underline C_{A|R}$ con valores en
$(A\text{-mod})$ es de \underline{tipo finito} si $F(B)$ es un
$A$-módulo de tipo finito siempre que $B$ sea una $R$-álgebra de tipo
finito. Por ejemplo, si $A$ es noetheriano y $M$ es un $A$-módulo de tipo
finito, entonces $\operatorname{Der}_R(-,M)$ es un haz aditivo de tipo
finito sobre $\underline C_{A|R}$.

Introducimos algo más de notación y recordamos algunos hechos básicos de
álgebra homológica.

\underline{3.1(a).} En todo el resto de esta sección,
$\widetilde{\underline C}_{A|R}$ y
$\overset{\vee}{\underline C}_{A|R}$ denotan, respectivamente, las
categorías de haces y prehaces con valores en $(A\text{-mod})$. Hay un
funtor de inclusión exacto por la izquierda
\[
 i_{A|R}:\widetilde{\underline C}_{A|R}
 \longrightarrow\overset{\vee}{\underline C}_{A|R}
\]
y un funtor exacto de haz asociado
\[
 G_{A|R}:\overset{\vee}{\underline C}_{A|R}
 \longrightarrow\widetilde{\underline C}_{A|R},
\]
adjunto por la izquierda de $i_{A|R}$.

%%%%%%%%%%  scan idx 83  |  folio 72  %%%%%%%%%%

Un objeto $B\in\operatorname{Ob}(\underline C_{A|R})$ induce un
funtor continuo
$j:\underline C_{B|R}\longrightarrow\underline C_{A|R}$ y, por tanto,
un diagrama conmutativo
\begin{equation}\tag{*}
\begin{tikzcd}[column sep=huge,row sep=huge]
 \widetilde{\underline C}_{A|R}
   \arrow[r,"\widetilde j"] \arrow[d,"i_{A|R}"']
 & \widetilde{\underline C}_{B|R} \arrow[d] \\
 \overset{\vee}{\underline C}_{A|R}
   \arrow[r,"j^{\vee}"']
 & \overset{\vee}{\underline C}_{B|R}.
\end{tikzcd}
\end{equation}
El funtor de restricción $\widetilde j$, al igual que $j^{\vee}$, es
exacto, porque el sitio $\underline C_{B|R}$ es simplemente la categoría
de las flechas de $\underline C_{A|R}$ con destino $B$, dotada de la
topología inducida.

Un homomorfismo de anillos $S\longrightarrow R$ induce un funtor continuo
$\beta:\underline C_{A|R}\longrightarrow\underline C_{A|S}$ y, por
tanto, un diagrama conmutativo
\begin{equation}\tag{**}
\begin{tikzcd}[column sep=huge,row sep=huge]
 \widetilde{\underline C}_{A|S}
   \arrow[r,"\widetilde\beta"] \arrow[d,"i_{A|S}"']
 & \widetilde{\underline C}_{A|R} \arrow[d] \\
 \overset{\vee}{\underline C}_{A|S}
   \arrow[r,"\beta^{\vee}"']
 & \overset{\vee}{\underline C}_{A|R}.
\end{tikzcd}
\end{equation}
Obsérvese que $\widetilde\beta$ no tiene por qué ser exacto.

\underline{3.1(b).} El funtor derivado por la derecha de la inclusión
\[
 i_{A|R}:\widetilde{\underline C}_{A|R}
 \longrightarrow\overset{\vee}{\underline C}_{A|R}
\]
se denotará por $\underline H^{\bullet}_{A|R}$. Para todo haz $F$ sobre
$\underline C_{A|R}$, pongamos
\[
 H^{\bullet}(A|R,F)=
 [\underline H^{\bullet}_{A|R}(F)](A).
\]

Para $B'\longrightarrow B$ en $\operatorname{Cov}(B)$, definamos
\[
 H^0(\{B'\longrightarrow B\},F)=
 \ker\!\left(
 F(B')\mathrel{\substack{\longrightarrow\\[-.6ex]\longrightarrow}}
 F(B'\mathbin{\times_B}B')
 \right).
\]
Los funtores derivados por la derecha de
$H^0(\{B'\longrightarrow B\},-)$ se denotan por
$H^{\bullet}(\{B'\longrightarrow B\},-)$. Estos son
los grupos de cohomología del módulo semisimplicial de \v{C}ech

%%%%%%%%%%  scan idx 84  |  folio 73  %%%%%%%%%%

\[
\begin{tikzcd}[column sep=large]
 F(B')
   \arrow[r,shift left=2] \arrow[r,shift right=2]
 & F(B'\times_BB')
   \arrow[r,shift left=3] \arrow[r] \arrow[r,shift right=3]
 & F(B'\times_BB'\times_BB')
   \arrow[r,shift left=4.5] \arrow[r,shift left=1.5]
   \arrow[r,shift right=1.5] \arrow[r,shift right=4.5]
 & \cdots .
\end{tikzcd}
\]

El funtor derivado por la derecha de
\[
 \check{\underline H}^{\,0}_{A|R}:
 \overset{\vee}{\underline C}_{A|R}
 \longrightarrow\overset{\vee}{\underline C}_{A|R}
\]
se denotará por $\check{\underline H}^{\,\bullet}_{A|R}$, donde
\[
 [\check{\underline H}^{\,0}_{A|R}(F)](B)=
 \varinjlim_{\{B'\to B\}\in\operatorname{Cov}(B)}
 H^0(\{B'\longrightarrow B\},F).
\]
Para todo prehaz $F$ sobre $\underline C_{A|R}$, pongamos
\[
 \check H^{\bullet}(A|R,F)=
 [\check{\underline H}^{\,\bullet}_{A|R}(F)](A)=
 \varinjlim_{\{A'\to A\}\in\operatorname{Cov}(A)}
 H^{\bullet}(\{A'\longrightarrow A\},F).
\]
Para todo haz $F$ sobre $\underline C_{A|R}$ se tiene la sucesión
espectral usual
\[
 \check H^p\bigl(A|R,\underline H^q_{A|R}(F)\bigr)
 \Longrightarrow H^{p+q}(A|R,F).
\]
Aquí
$\check H^0(A|R,\underline H^q_{A|R}(F))=0$ para todo $q>0$. En
consecuencia,
\[
 \check H^1(A|R,F)=H^1(A|R,F),
\]
y la sucesión
\[
 0\longrightarrow\check H^2(A|R,F)
 \longrightarrow H^2(A|R,F)
 \longrightarrow
 \check H^1\bigl(A|R,\underline H^1_{A|R}(F)\bigr)
 \longrightarrow\check H^3(A|R,F)
\]
es exacta.
